% !TeX root = ../../Book2.tex
% 纲要标题：### 4.1 Smith 标准形
% 纲要位置：计划纲要.md，第 915 行。

\section{Smith 标准形}
\label{b2:sec:0401}

% 纲要要点（供目录核对）：
% * 主理想整环上的幺模行列等价；Euclid 整环上的通常初等行列变换与算法
% * Smith 标准形与不变因子
% * 模展示与 Smith 标准形
% * 一般主理想整环上的存在证明使用 Bézout 型可逆变换和主理想的链条件，不预设任意可逆矩阵均可分解为通常三类初等矩阵
% * 分别陈述矩形矩阵、秩亏矩阵和零矩阵的标准形；不变因子的单位倍不确定性与归一化约定

本节设 \(R\) 为主理想整环，\(K\) 为其分式域。有限生成交换群中的整数关系矩阵曾把分类问题化为整除问题；这里把系数换成 \(R\)，并允许关系个数与生成元个数不同。首先要弄清可用的变换：在一般主理想整环上，存在可逆换基矩阵，不表示任意这样的矩阵都能由通常三类初等矩阵分解。

\subsection{幺模等价与初等行列算法}
\label{b2:subsec:0401:01}

采用第三章的列向量约定，对 \(A\in M_{m\times n}(R)\)，允许分别左乘 \(P\in\operatorname{GL}_m(R)\)、右乘 \(Q\in\operatorname{GL}_n(R)\)。若
\[
B=PAQ,
\]
称 \(A,B\) \term[yaomo-dengjia]{幺模等价}[unimodularly equivalent]。这里\term[yaomo-juzhen]{幺模矩阵}[unimodular matrix]指系数在 \(R\) 中、逆矩阵仍在 \(R\) 中的方阵；对交换环，它等价于行列式是单位。必要性来自行列式的乘法公式；充分性由伴随矩阵公式给出 \(P^{-1}=(\det P)^{-1}\operatorname{adj}(P)\)。

\ProseParagraphBreak
在展示中，左乘 \(P\) 改变目标自由模的生成元坐标，右乘 \(Q\) 重组来源自由模的关系坐标。若具体选定的新目标基、新来源基在旧基下的坐标矩阵分别是 \(U,V\)，同一映射的新矩阵仍是第三章的 \(U^{-1}AV\)：输入先由 \(V\) 换回旧坐标，输出再由 \(U^{-1}\) 换成新坐标。这里的 \(P\) 对应 \(U^{-1}\)，\(Q\) 对应 \(V\)；因逆矩阵也遍历全部可逆矩阵，幺模等价与这一换基公式一致。

\ProseParagraphBreak
交换两行，把一行乘一个单位，把一行加上另一行的任意倍，都是可逆行变换；列变换相同。逆操作分别为再次交换、乘逆单位与减去相同倍数。对 \(\mathbb Z\)，只能把行乘 \(1\) 或 \(-1\)，不能把一行随意除以非零整数。对 \(k[t]\)，可以乘任意非零常数，但不能把整行除以 \(t\)。

\ProseParagraphBreak
若 \(R\) 是 Euclid 整环，可以通过带余除法缩小一个非零主元。例如要清除主元 \(a\) 下方的 \(b\)，写 \(b=qa+r\)，先减去 \(q\) 倍主元行；若 \(r\ne0\)，交换两行，把主元换成 Euclid 函数值更小的 \(r\)。列方向采用相同做法。余数为零时直接消去该项。这个过程解释了整数与一元多项式上的实际算法；一般主理想整环未必配有 Euclid 函数，存在性证明须换用另一种终止理由。

\subsection{Smith 标准形与不变因子}
\label{b2:subsec:0401:02}

对角化使不同坐标上的关系分开，但这些循环关系还可能重新组合。例如整数矩阵 \(\operatorname{diag}(6,10)\) 已经对角，却有 \(6\nmid10\)；利用 Bézout 变换还可以把它化为 \(\operatorname{diag}(2,30)\)。要从这些不同的对角表达中确定标准形，还需要规定各非零对角元的整除次序。

\begin{definition}\label{b2:def:smith-normal-form}
设 \(R\) 是主理想整环，\(A\in M_{m\times n}(R)\)。若矩形对角矩阵 \(D\) 与 \(A\) 幺模等价，且具有形状
\[
D=\operatorname{diag}(d_1,\ldots,d_r,0,\ldots,0),
\qquad d_1\mid d_2\mid\cdots\mid d_r,
\]
其中 \(0\le r\le\min(m,n)\)，每个 \(d_i\)（\(1\le i\le r\)）均非零，所有非对角位置为零，就称 \(D\) 为 \(A\) 的\term[Smith-biaozhunxing]{Smith 标准形}[Smith normal form]。非零对角元称为矩阵的\term[bubian-yinzi]{不变因子}[invariant factors]。
\end{definition}

若把某个 \(d_i\) 乘上单位并相应调整对角形，所得仍是 Smith 标准形；因此唯一性首先是相伴类的唯一性。整数上通常取正数，\(k[t]\) 上取首一多项式；一般主理想整环没有默认的归一化。单位不变因子也属于固定尺寸矩阵的数据，不能直接删去。

\ProseParagraphBreak
各阶子式生成的理想将在\secref{b2:sec:0402}用来证明不变因子的唯一性。这里先看对角形怎样给出模的循环分解，再证明每个矩阵都能化为这种形状。

\subsection{模展示与 Smith 标准形}
\label{b2:subsec:0401:03}

\begin{proposition}\label{b2:prop:smith-cokernel}
设 \(R\) 是主理想整环，\(A\in M_{m\times n}(R)\)，\(P\in\operatorname{GL}_m(R)\)，\(Q\in\operatorname{GL}_n(R)\)。若 \(D=PAQ\) 是 \(A\) 的 Smith 标准形，非零对角元为 \(d_1,\ldots,d_r\)，则
\[
\operatorname{coker}A=R^m/A(R^n)
\cong\bigoplus_{i=1}^r R/(d_i)\ \oplus R^{m-r}.
\]
从原余核到对角余核的同构为 \(x+A(R^n)\mapsto Px+D(R^n)\)。
\end{proposition}
\begin{proof}
\(Q\) 满射，故 \(D(R^n)=PAQ(R^n)=P(A(R^n))\)。因此 \(P\) 把两个关系子模互相对应，诱导所写商模同构，逆由 \(P^{-1}\) 给出。

对角矩阵的像是
\[
d_1Re_1\oplus\cdots\oplus d_rRe_r\subseteq R^m.
\]
逐坐标把前 \(r\) 个坐标模去 \((d_i)\)，保留后 \(m-r\) 个坐标，得到从 \(R^m\) 到题述直和的满同态，其核恰是 \(D(R^n)\)。应用\cref{b2:thm:module-first-isomorphism}便得到结论。
\end{proof}

在这个分解中，非单位的非零因子 \(d_i\) 给出循环挠分量 \(R/(d_i)\)，单位因子则给出零模 \(R/(1)=0\)。这两类非零对角元都迫使相应的核坐标为零，因为整环中 \(d_ix_i=0\) 蕴含 \(x_i=0\)。零行是目标中没有受到关系约束的坐标，因而贡献余核的自由分量；零列是来源中映为零的坐标方向，因而贡献核中的自由方向。它们表示化简后的关系符号有冗余，并不增加余核的自由坐标。若需要在原生成元中指出某一循环分量的生成元，就要把对角坐标 \(e_i\) 沿商同构的逆映射送回，得到 \(P^{-1}e_i\) 的陪集；只记录对角元将丢失这项坐标信息。

\subsection{Bézout 变换与存在性证明}
\label{b2:subsec:0401:04}

主理想整环中，有限多个元素生成的理想仍有单个生成元。因此对非同时为零的 \(a,b\)，可以取 \(g\ne0\) 使 \((a,b)=(g)\)，并找到 \(u,v\in R\) 满足 \(ua+vb=g\)。我们希望把列向量 \((a,b)^{\mathsf T}\) 变为 \((g,0)^{\mathsf T}\)：第一行的系数已有 Bézout 等式给出的 \(u,v\)，第二行则要使两项相消，可以取 \(-b/g,a/g\)。这样得到
\[
E=\begin{pmatrix}u&v\\-b/g&a/g\end{pmatrix},
\qquad
E\begin{pmatrix}a\\b\end{pmatrix}=\begin{pmatrix}g\\0\end{pmatrix},
\qquad \det E=1.
\]
矩阵中的两个商属于 \(R\)，因为 \(g\) 同时整除 \(a,b\)；行列式为 \((ua+vb)/g=1\)，也正由同一 Bézout 等式保证。这是一项 \term[Bezout-bianhuan]{Bézout 变换}[Bézout transformation]。把它嵌入指定的两行得到幺模行变换；转置后右乘则处理同一行中的两项。这里直接给出了可逆矩阵，不需要把它拆成通常三类初等矩阵。

\begin{lemma}\label{b2:lem:pid-ascending-chain}
设 \(R\) 是主理想整环。\(R\) 的任意理想升链 \(I_1\subseteq I_2\subseteq\cdots\) 都最终稳定。
\end{lemma}
\begin{proof}
并集 \(I=\bigcup_n I_n\) 是理想：任意两个元素同属某个较大的 \(I_n\)，故它们的和、差与标量倍仍在并集中。写 \(I=(a)\)，则 \(a\in I_N\) 对某个 \(N\) 成立，所以 \(I\subseteq I_N\subseteq I\)。由此 \(I_n=I_N\) 对 \(n\ge N\) 成立。
\end{proof}

\begin{theorem}[Smith 标准形的存在性]\label{b2:thm:smith-normal-form}
设 \(R\) 是主理想整环，\(K\) 是 \(R\) 的分式域。对任意 \(A\in M_{m\times n}(R)\)，存在 \(P\in\operatorname{GL}_m(R)\)、\(Q\in\operatorname{GL}_n(R)\)，使 \(PAQ\) 为 Smith 标准形。其非零对角元个数等于 \(A\) 在 \(K\) 上的秩。
\end{theorem}
\begin{proof}
对 \(\min(m,n)\) 归纳。没有行或没有列时结论直接成立。零矩阵已经是标准形。否则通过交换行列把某个非零元素移到第一项，记当前主元为 \(a\ne0\)。为了递归处理其余行列，我们要隔离这个主元，并使它整除剩余块的全部条目；这样剩余块化成的对角元才能接在 \(a\) 后面形成整除链。

若第一列中有 \(b\) 不被 \(a\) 整除，对这两行施加上述 Bézout 变换，将主元换成 \(g\)，其中 \((g)=(a,b)\) 严格包含 \((a)\)。若第一行中有不被主元整除的项，对相应两列作转置的变换，得到同样的严格包含。一次这样的操作可能重新产生此前消去的项；重新检查第一行、第一列即可。

若第一行、第一列的所有项都被当前 \(a\) 整除，先用行减法清除第一列下方，再用列减法清除第一行右方。这只需有限步，而且第二轮不破坏已清空的第一列下方。此时得到
\[
\begin{pmatrix}a&0\\0&B\end{pmatrix}.
\]
若 \(a\) 不整除 \(B\) 的某项 \(b_{ij}\)，把含该项的行加到第一行。第一列主元仍为 \(a\)，第一行却出现不被 \(a\) 整除的 \(b_{ij}\)。再作一次列 Bézout 变换，主元理想又严格扩大，然后重新处理首行、首列。无论不整除的项来自首行、首列还是右下块，这些修正都使同一个主元理想严格增大。两次这样的增大之间，只需检查有限多个条目，并作有限次清行、清列操作；普通清零不改变主元。由\cref{b2:lem:pid-ascending-chain}，不整除修正不能发生无限多次，整个过程因而终止。最终得到上述分块矩阵，并且 \(a\) 整除 \(B\) 的每一项。

对 \(B\) 应用归纳假设。只在后面的行列内进行的幺模变换把 \(B\) 的新条目写成旧条目的 \(R\) 线性组合，所以它们仍被 \(a\) 整除。因此所得非零对角元满足 \(a\mid d_2\mid\cdots\mid d_r\)。把所有行列变换分别相乘，就得到所需 \(P,Q\)。每个变换都在 \(R\) 上可逆，在 \(K\) 上当然也可逆；对角形的 \(K\) 秩正是非零对角元数，所以等于原矩阵的秩。
\end{proof}

若 \(R\) 恰有 Euclid 除法，前述带余算法可以用 Euclid 函数的严格下降来实现这些约化，代替存在性证明中的理想升链论证。两种论证有相同的整除目标，却不预设相同的有效计算条件。

\subsection{矩形、秩亏与零矩阵的标准形}
\label{b2:subsec:0401:05}

对 \(m\times n\) 的秩为 \(r\) 的矩阵，Smith 标准形可以明确写成
\[
\begin{pmatrix}
\operatorname{diag}(d_1,\ldots,d_r)&0\\
0&0
\end{pmatrix},
\]
左上块为 \(r\times r\)，剩下的行数是 \(m-r\)，列数是 \(n-r\)。\(r=0\) 时整个矩阵为零；\(r=\min(m,n)\) 时也可能还有零行或零列，取决于哪一个尺寸更大。

\ProseParagraphBreak
例如
\[
D=\begin{pmatrix}1&0&0&0\\0&6&0&0\\0&0&0&0\end{pmatrix}
\]
是一个整数矩阵的合法标准形。它有两个非零不变因子，余核为 \(\mathbb Z/6\mathbb Z\oplus\mathbb Z\)，核为最后两个坐标组成的 \(\mathbb Z^2\)。一般地，\(\ker D\) 由最后 \(n-r\) 个标准向量生成。由 \(D=PAQ\) 且 \(P\) 可逆，
\[
Dy=0\quad\Longleftrightarrow\quad AQy=0.
\]
所以 \(Q\) 把 \(\ker D\) 映入 \(\ker A\)；反过来，对 \(x\in\ker A\) 取 \(y=Q^{-1}x\)，就有 \(Dy=PAx=0\)。因此 \(\ker A=Q(\ker D)\)，其基是最后 \(n-r\) 个向量 \(Qe_i\)。这与余核生成元使用 \(P^{-1}e_i\) 的方向不同：核向量要从新来源坐标送回原来源，余核生成元则要从新目标坐标送回原目标。

\ProseParagraphBreak
若 \(m=0\)，映射 \(R^n\to0\) 的余核为零；若 \(n=0\)，映射 \(0\to R^m\) 的余核为 \(R^m\)。这些情况同样符合公式，不需要另造“非零不变因子”。固定矩阵尺寸时保留单位对角元；转入模的同构分类时才删去它们所代表的零直和分量。

\ProseParagraphBreak
需要实际求出换基矩阵时，可继续阅读\secref{b2:sec:A01}的\cref{b2:alg:A-smith}。附录保留每一步行列操作，并在式\eqref{b2:eq:A-integer-certificate}与式\eqref{b2:eq:A-polynomial-certificate}中给出可直接相乘核验的证书。
